\documentclass[12pt]{article}

\input{../../_assets/preambulo.tex}


\begin{document}

    % 1. Foto de fondo
    % 2. Título
    % 3. Encabezado Izquierdo
    % 4. Color de fondo
    % 5. Coord x del titulo
    % 6. Coord y del titulo
    % 7. Fecha

    
    \input{../../_assets/portada}
    \portadaExamen{ffccA4.jpg}{Ecuaciones\\Diferenciales II\\Examen IX}{Ecuaciones Diferenciales II. Examen IX}{MidnightBlue}{-8}{28}{2026}{}

    \begin{description}
        \item[Asignatura] Ecuaciones Diferenciales II.
        \item[Curso Académico] 2023/2024.
        \item[Grado] Doble Grado en Ingeniería Informática y Matemáticas.
        \item[Grupo] Único.
        \item[Profesor] Rafael Ortega Ríos.
        \item[Descripción] Primera prueba de clase.
        \item[Fecha] 10 de Abril de 2024.
        % \item[Duración] ---.
    \end{description}
    \newpage


    % ------------------------------------
    
    \begin{ejercicio}
        Nos dan dos funciones de clase $C^1$,
        \Func{X_i}{\bb{R} \times \bb{R}}{\bb{R}}{(t,x)}{X_i(t,x), \quad i = 1,2,}
        y se supone que $X_1(t,x) = X_2(t,x)$ si $t \in \bb{R}$, $|x| \leq 2$. Para cada $i = 1,2$ la función $x_i : I_i \longrightarrow \bb{R}$, $t \longmapsto x_i(t)$ es una solución del problema
        \[
        x' = X_i(t,x), \quad x(0) = 0.
        \]
        ¿Se puede asegurar que $x_1(t)$ y $x_2(t)$ coinciden en un entorno de $t_0 = 0$?
    \end{ejercicio}

    \begin{ejercicio}
        Demuestra la existencia de un número real $r > 0$ y dos funciones continuas $x, y : \left[-r,r\right] \longrightarrow \bb{R}$ que cumplen
        \[
        x(t) = 2e^{\int\limits_0^t y(s)ds}, \quad y(t) = e^{-3\int\limits_0^t x(s)ds},
        \]
        para cada $t \in \left[-r,r\right]$.
    \end{ejercicio}

    \begin{ejercicio}
        Encuentra una solución maximal del problema
        \[
        x' = t^2\sqrt{x^+}, \quad x(0) = 1
        \]
        donde $x^+ = \max(x,0)$. ¿Hay unicidad local? ¿y global?
    \end{ejercicio}

    \begin{ejercicio}
        Dado el intervalo abierto $I = \left]-1,1\right[$ se supone que $\cc{F} \subset C(I)$ es un conjunto de funciones continuas $f : I \longrightarrow \bb{R}$ con la siguiente propiedad:
        \begin{center}
            \emph{Si $f_1, f_2 \in \cc{F}$ cumplen $f_1(t_0) = f_2(t_0)$ para algún $t_0 \in I$, entonces existe $\delta > 0$ tal que $f_1(t) = f_2(t)$, $\forall t \in \left]t_0 - \delta, t_0 + \delta\right[$.}
        \end{center}

        Demuestra que existe a lo sumo una función que cumple $f \in \cc{F}$, $f(0) = 3$.
    \end{ejercicio}

    \begin{ejercicio}
        Se supone que $\{x_n\}_{n \geq 1}$ es una sucesión de funciones en $C^1\left[-1,1\right]$ que cumplen, para cada $t \in \left[-1,1\right]$,
        \[
        |\dot{x}_n(t)| \leq 6, \quad \left|x_n(t) - 7 - 3\int_0^t \frac{x_n(s)}{1+x_n(s)^2} ds\right| \leq \frac{1}{n}.
        \]
        \begin{enumerate}[label=\alph*)]
            \item Prueba que existe una sucesión parcial de $\{x_n\}$ que converge uniformemente en $\left[-1,1\right]$ a una función límite $x(t)$.
            \item Encuentra un problema de valores iniciales del que es solución $x(t)$.
            \item Prueba que la sucesión $\{x_n\}$ converge uniformemente.
        \end{enumerate}
    \end{ejercicio}

    \newpage
    \setcounter{ejercicio}{0} % Reiniciar contador de ejercicios
    \noindent
    % \textbf{Solución.}

\end{document}